% Source: precision_vs_recall.tex
\begin{notechapter}
\chapter{Class Encoding, Precision and Recall}\label{ch:precision-recall}

\section{Introduction}
For binary classification, the importance of encoding the target class
as the positive (encoded as 1) or negative (encoded as 0)
depends on whether the classes are 
well-balanced or unbalanced. Let's explore two scenarios with 
examples and computations:
\begin{itemize}
    \item \textbf{Scenario 1}: A classifier with default settings.
    \item \textbf{Scenario 2}: A classifier with a modified threshold 
    to better identify harmful cases, at the cost of more 
    misclassifications for non-harmful cases.
\end{itemize}

Each scenario will be analyzed for both balanced and unbalanced 
datasets, considering two encoding choices:
\begin{itemize}
    \item Encoding the harmful class (poisonous mushrooms) as 1.
    \item Encoding the harmful class as 0.
\end{itemize}

This will highlight how encoding affects precision and recall in 
different contexts.

\section{Scenario 1: Default Classifier}
In this scenario, the classifier has default settings, resulting 
in higher misclassification rates for the harmful class.

\subsection{Balanced Dataset}
\begin{itemize}
    \item Dataset: 50,000 edible mushrooms and 50,000 poisonous 
    mushrooms.
    \item Misclassification rates:
        \begin{itemize}
            \item 20\% of poisonous mushrooms.
            \item 5\% of edible mushrooms.
        \end{itemize}
\end{itemize}

\subsubsection{Encoding Poisonous as 1}
\begin{itemize}
    \item Confusion Matrix:
    \[
    \begin{array}{c|cc}
        & \text{Predicted Poisonous} & \text{Predicted Edible} \\
        \hline
        \text{Actual Poisonous} & 40,000 & 10,000 \\
        \text{Actual Edible} & 2,500 & 47,500 \\
    \end{array}
    \]
    \item Recall (for poisonous mushrooms):
    \[
    \text{Recall} = \frac{40000}{40000 + 10000} = 0.8
    \]
    \item Precision (for poisonous mushrooms):
    \[
    \text{Precision} = \frac{40000}{40000 + 2500} \approx 0.941
    \]
\end{itemize}

\subsubsection{Encoding Poisonous as 0}
\begin{itemize}
    \item Confusion Matrix:
    \[
    \begin{array}{c|cc}
        & \text{Predicted Edible} & \text{Predicted Poisonous} \\
        \hline
        \text{Actual Edible} & 47,500 & 2,500 \\
        \text{Actual Poisonous} & 10,000 & 40,000 \\
    \end{array}
    \]
    \item Recall (for edible mushrooms):
    \[
    \text{Recall} = \frac{47500}{47500 + 2500} = 0.95
    \]
    \item Precision (for edible mushrooms):
    \[
    \text{Precision} = \frac{47500}{47500 + 10000} \approx 0.826
    \]
\end{itemize}

\subsection{Unbalanced Dataset}
\begin{itemize}
    \item Dataset: 100,000 edible mushrooms and 1,000 poisonous 
    mushrooms.
    \item Misclassification rates:
        \begin{itemize}
            \item 20\% of poisonous mushrooms.
            \item 5\% of edible mushrooms.
        \end{itemize}
\end{itemize}

\subsubsection{Encoding Poisonous as 1}
\begin{itemize}
    \item Confusion Matrix:
    \[
    \begin{array}{c|cc}
        & \text{Predicted Poisonous} & \text{Predicted Edible} \\
        \hline
        \text{Actual Poisonous} & 800 & 200 \\
        \text{Actual Edible} & 5,000 & 95,000 \\
    \end{array}
    \]
    \item Recall (for poisonous mushrooms):
    \[
    \text{Recall} = \frac{800}{800 + 200} = 0.8
    \]
    \item Precision (for poisonous mushrooms):
    \[
    \text{Precision} = \frac{800}{800 + 5000} \approx 0.138
    \]
\end{itemize}

\subsubsection{Encoding Poisonous as 0}
\begin{itemize}
    \item Confusion Matrix:
    \[
    \begin{array}{c|cc}
        & \text{Predicted Edible} & \text{Predicted Poisonous} \\
        \hline
        \text{Actual Edible} & 95,000 & 5,000 \\
        \text{Actual Poisonous} & 200 & 800 \\
    \end{array}
    \]
    \item Precision (for edible mushrooms):
    \[
    \text{Precision} = \frac{95000}{95000 + 200} \approx 0.998
    \]
    \item Recall (for edible mushrooms):
    \[
    \text{Recall} = \frac{95000}{95000 + 5000} = 0.95
    \]
\end{itemize}

\subsection{Key Observations for Scenario 1}
\begin{itemize}
    \item In the balanced case, encoding the harmful class as 1 
    emphasizes recall (0.8) and precision (0.941) for poisonous 
    mushrooms. Encoding it as 0 emphasizes recall (0.95) and 
    precision (0.826) for edible mushrooms. Both metrics are 
    meaningful and comparable.
    \item In the unbalanced case, encoding the harmful class as 1 
    highlights the recall metric (0.8), which is critical for 
    identifying harmful cases. However, the precision is very low 
    (0.138), reflecting the high cost of false positives.
    \item Encoding the harmful class as 0 masks the severity of 
    the problem, as the precision and recall for the majority class 
    are very close to 1, making the classifier appear to perform 
    well despite failing to identify 20\% of harmful cases.
\end{itemize}

\section{Scenario 2: Modified Threshold}
In this scenario, the classifier has a modified threshold to 
better identify harmful cases, resulting in fewer misclassifications 
for the harmful class but more misclassifications for the non-harmful 
class.

\subsection{Balanced Dataset}
\begin{itemize}
    \item Dataset: 50,000 edible mushrooms and 50,000 poisonous 
    mushrooms.
    \item Misclassification rates:
        \begin{itemize}
            \item 1\% of poisonous mushrooms.
            \item 10\% of edible mushrooms.
        \end{itemize}
\end{itemize}

\subsubsection{Encoding Poisonous as 1}
\begin{itemize}
    \item Confusion Matrix:
    \[
    \begin{array}{c|cc}
        & \text{Predicted Poisonous} & \text{Predicted Edible} \\
        \hline
        \text{Actual Poisonous} & 49,500 & 500 \\
        \text{Actual Edible} & 5,000 & 45,000 \\
    \end{array}
    \]
    \item Recall (for poisonous mushrooms):
    \[
    \text{Recall} = \frac{49500}{49500 + 500} \approx 0.990
    \]
    \item Precision (for poisonous mushrooms):
    \[
    \text{Precision} = \frac{49500}{49500 + 5000} \approx 0.908
    \]
\end{itemize}

\subsubsection{Encoding Poisonous as 0}
\begin{itemize}
    \item Confusion Matrix:
    \[
    \begin{array}{c|cc}
        & \text{Predicted Edible} & \text{Predicted Poisonous} \\
        \hline
        \text{Actual Edible} & 45,000 & 5,000 \\
        \text{Actual Poisonous} & 500 & 49,500 \\
    \end{array}
    \]
    \item Recall (for edible mushrooms):
    \[
    \text{Recall} = \frac{45000}{45000 + 5000} = 0.9
    \]
    \item Precision (for edible mushrooms):
    \[
    \text{Precision} = \frac{45000}{45000 + 500} \approx 0.989
    \]
\end{itemize}

\subsection{Unbalanced Dataset}
\begin{itemize}
    \item Dataset: 100,000 edible mushrooms and 1,000 poisonous 
    mushrooms.
    \item Misclassification rates:
        \begin{itemize}
            \item 1\% of poisonous mushrooms.
            \item 10\% of edible mushrooms.
        \end{itemize}
\end{itemize}

\subsubsection{Encoding Poisonous as 1}
\begin{itemize}
    \item Confusion Matrix:
    \[
    \begin{array}{c|cc}
        & \text{Predicted Poisonous} & \text{Predicted Edible} \\
        \hline
        \text{Actual Poisonous} & 990 & 10 \\
        \text{Actual Edible} & 10,000 & 90,000 \\
    \end{array}
    \]
    \item Recall (for poisonous mushrooms):
    \[
    \text{Recall} = \frac{990}{990 + 10} = 0.990
    \]
    \item Precision (for poisonous mushrooms):
    \[
    \text{Precision} = \frac{990}{990 + 10000} \approx 0.090
    \]
\end{itemize}

\subsubsection{Encoding Poisonous as 0}
\begin{itemize}
    \item Confusion Matrix:
    \[
    \begin{array}{c|cc}
        & \text{Predicted Edible} & \text{Predicted Poisonous} \\
        \hline
        \text{Actual Edible} & 90,000 & 10,000 \\
        \text{Actual Poisonous} & 10 & 990 \\
    \end{array}
    \]
    \item Precision (for edible mushrooms):
    \[
    \text{Precision} = \frac{90000}{90000 + 10} \approx 0.9999
    \]
    \item Recall (for edible mushrooms):
    \[
    \text{Recall} = \frac{90000}{90000 + 10000} = 0.9
    \]
\end{itemize}

\subsection{Key Observations for Scenario 2}
\begin{itemize}
    \item In the balanced case, encoding the harmful class as 1 
    emphasizes recall (0.990) and precision (0.908) for poisonous 
    mushrooms. Encoding it as 0 emphasizes recall (0.9) and 
    precision (0.989) for edible mushrooms. Both metrics are 
    meaningful and comparable.
    \item In the unbalanced case, encoding the harmful class as 1 
    highlights the recall metric (0.990), which is critical for 
    identifying harmful cases. However, the precision is very low 
    (0.090), reflecting the high cost of false positives.
    \item Encoding the harmful class as 0 masks the severity of 
    the problem, as the precision and recall for the majority class 
    are very close to 1, making the classifier appear to perform 
    well despite failing to identify 1\% of harmful cases.
\end{itemize}

\section{Conclusion}
\begin{itemize}
    \item For \textbf{balanced classes}, the encoding choice 
    affects which class's precision and recall you focus on, 
    but both metrics are similarly expressive and meaningful.
    \item For \textbf{unbalanced classes}, encode the minority 
    class as 1 and prioritize recall, especially when 
    misclassifying the minority class is harmful.
    \item The computations and confusion matrices clearly show 
    why encoding the harmful class as 1 is crucial in unbalanced 
    cases, as it ensures that the metrics reflect the performance 
    on the critical class.
    \item Comparing Scenario 1 and Scenario 2 highlights how 
    optimizing the classifier to reduce misclassifications for 
    the harmful class improves recall but often at the cost of 
    lower precision.
\end{itemize}
\end{notechapter}
